\usepackage[margin=1in]{geometry}
\usepackage{libertine}
\usepackage[T1]{fontenc}
\usepackage[utf8]{inputenc}
\usepackage{amsfonts,amssymb,amsmath,amsthm}
% \usepackage[sort&compress]{natbib}
%\usepackage[english]{babel}
%\usepackage{eurosym}
\usepackage{calc}
\usepackage{tikz}
\usetikzlibrary{calc}

%\usepackage{verbatim,refcount}
%\usepackage{fancybox,soul,color}
\usepackage{url}
\usepackage{hyperref}
%\usepackage{subfig}
\usepackage{enumitem}
%\usepackage{epsfig}
%\usepackage[active]{srcltx}
%\usepackage[mathlines]{lineno}
\usepackage{xspace}
%\usepackage{multirow,tabu}
%\usepackage{slashed}
\usepackage{todonotes}
%\usepackage{eucal}
%\usepackage{mathtools}
\usepackage{thmtools}
\usepackage{mathrsfs}  
\usepackage{textpos}
\usepackage[normalem]{ulem}
\usepackage[capitalize, nameinlink]{cleveref}

%\usetikzlibrary{arrows,shapes,automata,backgrounds,petri, fit, decorations.pathreplacing,calligraphy}
%\DeclareMathOperator{\mad}{mad}
%\DeclareMathOperator{\ad}{ad}
%\DeclareMathOperator{\ntn}{ntn}
%\DeclareMathOperator{\func}{fun}
%\DeclareMathOperator{\degen}{degeneracy}
%\DeclareMathOperator{\avgd}{avgdeg}
%\DeclareMathOperator{\cwd}{cw}
%\DeclareMathOperator{\twin}{tww}
%\DeclareMathOperator{\fgn}{fgn}
\DeclareMathOperator{\VC}{2VCdim}
\DeclareMathOperator{\tw}{tw}
\DeclareMathOperator{\bsn}{bsn}
\DeclareMathOperator{\tpw}{tpw}
\DeclareMathOperator{\td}{td}
\DeclareMathOperator{\cc}{cc}
\DeclareMathOperator{\ddim}{ddim}

\newtheorem{problem}{Problem}[section]
\newtheorem{theorem}{Theorem}[section]
\newtheorem{lemma}[theorem]{Lemma}
\newtheorem{claim}[theorem]{Claim}
\newtheorem{corollary}[theorem]{Corollary}
\newtheorem{conjecture}[theorem]{Conjecture}
\theoremstyle{definition}
\newtheorem{definition}[theorem]{Definition}
\newcommand{\Oh}{\mathcal{O}}
\renewcommand{\varnothing}{\emptyset}

\def\Cc{\mathscr{C}}

\def\R{\mathbb{R}}
\def\C{\mathcal{C}}
\def\H{\mathcal{H}}
\def\P{\mathcal{P}}
\def\F{\mathcal{F}}
\def\N{\mathbb{N}}
\def\T{\mathcal{T}}
\def\D{\mathcal{D}}
\def\B{\mathrm{Ball}}
\def\Sol{\texttt{Sol}}
\newcommand{\Co}{\mathsf{Cont}}
\newcommand{\CMSO}{\textsf{CMSO}$_2$\xspace}
\newcommand{\cmsotwo}{\CMSO}
\newcommand{\VV}{$\bigvee\mkern-15mu\bigvee$}
\newcommand{\weight}{\mathsf{w}}

\newcommand{\Aa}{\mathcal{A}}
\newcommand{\Tt}{\mathcal{T}}
\newcommand{\Bb}{\mathcal{B}}
\newcommand{\Ss}{\mathcal{S}}
\newcommand{\Dd}{\mathcal{D}}
\newcommand{\Ww}{\mathcal{W}}

\newcommand{\wh}[1]{\widehat{#1}}

\newcommand{\bag}{\mathsf{bag}}
\newcommand{\dist}{\mathrm{dist}}
\newcommand{\rad}{\mathrm{rad}}

\newcommand{\Jadwiga}[2][]{\todo[color=yellow!80,inline,#1]{{\textbf{J:} #2}}}
\newcommand{\kk}[2][]{\todo[color=green!80,inline,#1]{{\textbf{kk:} #2}}}
\newcommand{\prz}[2][]{\todo[color=olive!80,#1]{{\textbf{prz:} #2}}}
\newcommand{\przin}[2][]{\todo[color=olive!80,#1,inline]{{\textbf{prz:} #2}}}
\newcommand{\mipilin}[2][]{\todo[color=blue!30,inline,#1]{{\textbf{MiPi:} #2}}}
\newcommand{\mipil}[2][]{\todo[color=blue!30,#1]{{\textbf{MiPi:} #2}}}
\newcommand{\mpil}[2][]{\todo[color=magenta!80,#1]{{\textbf{MP:} #2}}}

\newcommand{\tara}[2][]{\todo[color=violet!30,inline,#1]{{\textbf{Tara:} #2}}}

\newcommand{\ceil}[1]{\left\lceil #1 \right\rceil }

\renewcommand{\leq}{\leqslant}
\renewcommand{\geq}{\geqslant}
\renewcommand{\le}{\leqslant}
\renewcommand{\ge}{\geqslant}
\renewcommand{\mid}{~|~}

%\linenumbers

%%%%%%%%%%%%%%%%%%%%%%%%%%%%

%%%%%%%%%%%%%%%%%%%%%%%%%%%%%%%%%%%%%%%%%%%%%%%%%%%

\usepackage{framed}
\usepackage{tabularx}

\newlength{\RoundedBoxWidth}
\newsavebox{\GrayRoundedBox}
\newenvironment{GrayBox}[1]%
   {\setlength{\RoundedBoxWidth}{.93\columnwidth}
    \def\boxheading{#1}
    \begin{lrbox}{\GrayRoundedBox}
       \begin{minipage}{\RoundedBoxWidth}}%
   {   \end{minipage}
    \end{lrbox}
    \begin{center}
    \begin{tikzpicture}%
       \node(Text)[draw=black!20,fill=white,rounded corners,inner sep=2ex,text width=\RoundedBoxWidth]
             {\usebox{\GrayRoundedBox}};
        \coordinate(x) at (current bounding box.north west);
        \node [draw=white,rectangle,inner sep=3pt,anchor=north west,fill=white]
        at ($(x)+(6pt,.75em)$) {\boxheading};
    \end{tikzpicture}
    \end{center}}

\newenvironment{defproblemx}[1]{\noindent\ignorespaces%
                                \FrameSep=6pt%
                                \parindent=0pt%
%                 \vspace*{-1em}
                \begin{GrayBox}{#1}%
                \begin{tabular*}{\columnwidth}{!{\extracolsep{\fill}}@{\hspace{.1em}} >{\itshape} p{1.5cm} p{0.86\columnwidth} @{}}%
            }{
                \end{tabular*}%
                \end{GrayBox}%
                \ignorespacesafterend
%                 \vspace*{-1em}
            }

\newcommand{\problemTask}[3]{%
  \begin{defproblemx}{#1}
    Input: & #2 \\
    Task: & #3
  \end{defproblemx}
}

\newcommand{\parProblemTask}[4]{%
  \begin{defproblemx}{#1}
    Input:  & #2 \\
    Par.: & #3 \\
    Task: & #4
  \end{defproblemx}
}

\newcommand{\problemQuestion}[3]{%
  \begin{defproblemx}{#1}
    Input: & #2 \\
    Question: & #3
  \end{defproblemx}
}

\newcommand{\parProblemQuestion}[4]{%
  \begin{defproblemx}{#1}
    Input:  & #2 \\
    Par.: & #3 \\
    Question: & #4
  \end{defproblemx}
}

%\newcommand{\tara}[1]{{\color{violet}{(T: #1)}}}
\usepackage{mathtools}

\newcommand{\dd}{\hbox{-}}


\newenvironment{claimproof}[1][Proof of Claim.]{%
  \renewcommand{\qedsymbol}{$\lrcorner$}%
  \begin{proof}[#1]%
}{%
  \end{proof}%
}